\begin{frame}{Decision Variable}

\justifying

\begin{exampleblock}{$g_\mathrm{lim}$: repair threshold (mm)}
\justifying
Distance from the carbonation front to the rebar that triggers a repair.
It defines the whole maintenance schedule $t_1, t_2, \dots$
\end{exampleblock}

\vspace{0.3cm}

\begin{columns}[T,onlytextwidth]
  \column{0.48\textwidth}
  \textbf{Large $g_\mathrm{lim}$}
  \begin{itemize}
    \item Repair early
    \item Many repairs
    \item Few failures
  \end{itemize}

  \column{0.48\textwidth}
  \textbf{Small $g_\mathrm{lim}$}
  \begin{itemize}
    \item Repair late
    \item Few repairs
    \item Many failures
  \end{itemize}
\end{columns}

\vspace{0.3cm}

\begin{block}{Why a threshold and not a fixed interval?}
\justifying
The original concrete and the repair mortar carbonate at different rates.
A threshold in mm gives a different interval for each cycle, and it is the language engineers use in the field.
\end{block}

\end{frame}
